\documentclass{beamer}

\usepackage{biblatex}
\addbibresource{smc.bib}
\usepackage{hyperref}
\usepackage[utf8]{inputenc}

\usetheme{Copenhagen}
\usecolortheme{default}
\setbeamertemplate{items}[triangle]
\setbeamertemplate{sections/subsections in toc}[square]

%Information to be included in the title page:
\title{Bayesian parameter synthesis for markov population model.}
\author{Nhat-Huy Phung}
\institute{University of Konstanz}
\titlegraphic{\includegraphics[width=5cm]{figures/unisignet-klein.jpg}}


\begin{document}
\frame{\titlepage}

\begin{frame}
    \begin{center}
        \Huge Bayesian framework for parameter synthesis.
    \end{center}
\end{frame}

\begin{frame}
    \frametitle{Framework}
    %% General framework
    \begin{block}{Definition}
        Model checking is an automated technique that, given a finite-state model of a system and a formal property,systematically checks whether this property holds for (a given state in) that model
    \end{block}

\end{frame}


\begin{frame}
    \frametitle{Model checking step}
    Properties are specified by temporal logics
    Since we are interested in probabilistic model checking, our models encompasses probabilistic behaviours.
    \begin{examples}[Probabilistic models]
        \begin{itemize}
            \item Rational fucntion evaluation
            \item Statistical model checking
        \end{itemize}
    \end{examples}
\end{frame}

\begin{frame}
    \frametitle{Statistical model checking}
    Statistical Model Checking (SMC) is a formal verification technique that combines simulation and statistical methods for the analysis of stochastic systems. \footnotemark[1] Statistical Model Checking verifies a system $S$ property $\phi$ over a finite set of \textit{traces}, acquired through simulating the system of concern $S$.
    \begin{block}{Advantages}
        \begin{itemize}
            \item \textbf{Scalability}: avoid state space explosion issues.
        \end{itemize}
    \end{block}
    \footnotetext[1]{\href{https://www-verimag.imag.fr/Statistical-Model-Checking-814.html}{}}
\end{frame}


\begin{frame}
    \frametitle{Statistical model checking}
    Given a model $M$ of a system $S$ and a temporal property $\phi$. Let $p:=Pr\{M \models \phi \}$ be the probability that the model $M$ satisfies the property $\phi$. \footnotemark[1]
    \begin{block}{Verification}
        \begin{itemize}
            \item \textbf{Quantitative}: Estimate $p$
            \item \textbf{Qualitative}: Given a threshold $\theta$, test the hypothesis $H := $
        \end{itemize}
    \end{block}
    \footnotetext[1]{\href{https://www-verimag.imag.fr/Statistical-Model-Checking-814.html}{}}
\end{frame}

\begin{frame}
    \frametitle{Statistical model checking}
    \begin{block}{Quantitative}
        Estimate $p:=Pr\{M \models \phi \}$ wrt. precision $\delta$ and confidence level $\alpha$
    \end{block}
    The estimation is described in detail in \cite{agha2018survey}. We calculate $\hat{p}$ as an estimation of $p$ such that
    \begin{align*}
        Pr\{ |p - \hat{p}| < \delta \} = 1 - \alpha
    \end{align*}
    $\hat{p}$ can be estimated using different bounds, such as Chernoff-Hoeffding bound \cite{hoeffding1963probability}, Okamoto bound \cite{okamoto1959some} or Massart bound \cite{massart1990tight}
\end{frame}

\begin{frame}
    \frametitle{Statistical model checking}
    \begin{block}{Quantitative}
        Let $p:=Pr\{M \models \phi \}$ and a threshold $\theta \in [0,1]$. Compare $p$ and $\theta$
    \end{block}
    The general approach is to do hypothesis test \cite{younes2005verification} given a confidential level $\alpha$
    \begin{itemize}
        \item $H_0: p >= \theta$
        \item $H_1: p < \theta$
    \end{itemize}
    More details can be found at \cite{wald1945sequential}
\end{frame}

\begin{frame}
    \frametitle{Case study}
    PRISM \cite{kwiatkowska2002prism} is a model checking tool that support discrete event simulation and statistical model checking.
    For an example of how to use PRISM for statistical model checking please follow \href{https://www.prismmodelchecker.org/manual/RunningPRISM/StatisticalModelChecking}{this link}.
\end{frame}

\begin{frame}[allowframebreaks]{References}
    \printbibliography
\end{frame}

\end{document}